% !TEX root =  master.tex

\chapter{Technologien}

\paragraph{FFMPEG}  \label{par:ffmpeg}
FFMPEG bietet auf der einen Seite eine Abstraktionsebene über den Schnittstellen zu Video und Audio Geräten. Und auf der anderen Seite bietet es ein umfangreiches Sortiment an Funktionen, um Video- und Audiodaten zu bearbeiten.
es wird im Laufe dieser Arbeit verwendet, um Kameras in den HB10 zu integrieren, und Anschließend entstehende Videodaten zu komprimieren und für das Streaming mittels HLS (s. \ref{par:hls}) zu bearbeiten.

\paragraph{H.264/H.265 Codecs} \label{par:codecs}
H.264 und H.265 sind zwei Video-Codecs, welche eine Verlustbehaftete Codierung von Videos, mit intra- und interframe Komprimierung ermöglichen.
Diese Codecs werden als Kandidaten für die Komprimierung der Videos einbezogen.

\paragraph{POSIX} \label{par:posix}
POSIX ist der Standard für das lesen und schreiben von Dateien unter Linux. POSIX wird verwendet, um die Kommunikation zweier Apps auf dem HB10 über eine Socket zu implementieren.

%\paragraph{C++}

\paragraph{REST} \label{par:rest}
REST ist ein Regelwerk für http Zugriffe über das Internet. In dieser Arbeit wird es für die Übertragung der Videodaten zwischen sowohl Server und HB10, als auch zwischen Server und Nutzer verwendet.

\paragraph{HLS} \label{par:hls}
HLS ist ein HTTP-basiertes Video Streaming Protokoll bei dem der Client den Server nach einer playlist(.m3u8) Datei fragt welche alle Notwendigen Videodateien und deren Länge enthält, sodass der Client nicht das gesamte Video auf einmal herunterladen muss.
HLS wird verwendet, um die Videos vom Server an den Nutzer zu Streamen.

\paragraph{Spring Boot} \label{par:spring_boot}

Spring Boot ist eine Erweiterung des Spring Frameworks, und ermöglicht es einfach einen Webserver oder REST APIs zu entwickeln. Spring Boot wird verwendet um die Schnittstellen des Servers zu erstellen.

\paragraph{PostgreSql} \label{par:postgre}
PostgreSql ist ein Datenbank Framework, und erlaubt es einfach Datenbanken aufzusetzen. Es wird für eine relationale Datenbank verwendet, welche die Metadaten zu den Videos verwaltet.

\newpage